\documentclass[11pt,a4paper]{article}

\usepackage{amsmath,amssymb,amsthm}
\usepackage{algorithm}
\usepackage{algpseudocode}
\usepackage{graphicx}
\usepackage{hyperref}
\usepackage{booktabs}
\usepackage{xcolor}
\usepackage{geometry}
\geometry{margin=2.5cm}

\newtheorem{definition}{Definition}

\title{\textbf{Methodology of the \texttt{rna\_design} Project:\\
  EDA-Based RNA Inverse Folding with Structural Priors}}
\author{Roberto Santana\\
  Intelligent Systems Group, University of the Basque Country (UPV/EHU)\\
  \texttt{roberto.santana@ehu.eus}}
\date{}

\begin{document}
\maketitle
\tableofcontents
\bigskip

% ─────────────────────────────────────────────────────────────────────────────
\section{Problem Definition}
\label{sec:problem}
% ─────────────────────────────────────────────────────────────────────────────

Ribonucleic acid (RNA) molecules are single-stranded polymers composed of
four nucleotide bases: adenine (A), cytosine (C), guanine (G), and uracil
(U). Unlike DNA, RNA folds back on itself to form intricate secondary and
tertiary structures that are directly responsible for its biological
function. The secondary structure — the set of intramolecular base pairs
that form through Watson--Crick (A–U, G–C) and wobble (G–U) hydrogen bonds
— plays a central role in RNA catalysis, gene regulation, viral replication,
and protein translation~\cite{tinoco1999rna}.

The canonical secondary structure of an RNA sequence of length $n$ can be
described compactly in \emph{Dot-Bracket Notation} (DBN): a string $s$ of
length $n$ over the alphabet $\{$\texttt{(}, \texttt{)}, \texttt{.}$\}$, where
\texttt{(} and \texttt{)} mark the left and right sides of a Watson--Crick pair,
respectively, and \texttt{.} denotes an unpaired position. For example:
\begin{center}
  \texttt{Sequence:    A\;C\;G\;U\;A\;C\;G\;U}\\
  \texttt{Structure:   (\;(\;.\;.\;.\;.\;)\;)}
\end{center}
indicates that positions 1 and 8 form one base pair and positions 2 and 7
form another, while positions 3–6 are unpaired.

\subsection{RNA Inverse Folding}

Standard RNA \emph{forward} folding takes a nucleotide sequence as input and
predicts its minimum free-energy (MFE) secondary structure using
thermodynamic models~\cite{zuker1989mfold, hofacker1994fast}. The
\emph{inverse} problem — also called \textbf{RNA design} or
\textbf{RNA inverse folding} — is the reverse: given a target secondary
structure $s^*$, find one or more nucleotide sequences $\mathbf{x}$ such
that the MFE structure predicted for $\mathbf{x}$ is $s^*$ (or as close to
it as possible).

\begin{definition}[RNA Inverse Folding]
Let $s^* \in \{\texttt{(},\texttt{)},\texttt{.}\}^n$ be a target secondary
structure of length $n$. The RNA inverse folding problem is to find a
sequence $\mathbf{x} \in \{\mathrm{A,C,G,U}\}^n$ such that
\[
  \mathcal{F}(\mathbf{x}) = s^*,
\]
where $\mathcal{F}(\mathbf{x})$ denotes the MFE secondary structure
predicted for sequence $\mathbf{x}$ by a thermodynamic folding algorithm.
\end{definition}

RNA inverse folding is computationally hard in the general case. The solution
space is $4^n$, and the fitness landscape is highly rugged and deceptive: many
sequences may fold into very similar structures, while small single-nucleotide
changes can drastically alter the predicted structure. The problem is
therefore well-suited to population-based metaheuristics that can exploit
structural regularities.

\subsection{Relevance}

The ability to design RNA sequences with prescribed structures has broad
practical importance:
\begin{itemize}
  \item \textbf{Synthetic biology}: engineering functional RNA molecules such
    as ribozymes, riboswitches, aptamers, and siRNA for therapeutic
    applications~\cite{wachsmuth2013riboswitch};
  \item \textbf{Vaccine development}: mRNA sequences must fold into stable
    structures to ensure efficient translation and immunogenicity;
  \item \textbf{Structural biology}: generating synthetic sequences for
    experimental validation of predicted structures;
  \item \textbf{Citizen science}: the EteRNA project demonstrates that human
    puzzle solvers and algorithms can jointly advance our understanding of
    RNA folding~\cite{lee2014eterna}.
\end{itemize}

% ─────────────────────────────────────────────────────────────────────────────
\section{Problem Representation and Objective Functions}
\label{sec:representation}
% ─────────────────────────────────────────────────────────────────────────────

\subsection{Solution Representation}

Each candidate solution is a discrete vector
$\mathbf{x} = (x_1, x_2, \ldots, x_n) \in \{0,1,2,3\}^n$,
where each component $x_i$ encodes the nucleotide at position $i$ of the
RNA sequence using the mapping
\[
  0 \leftrightarrow \mathrm{A}, \quad
  1 \leftrightarrow \mathrm{C}, \quad
  2 \leftrightarrow \mathrm{G}, \quad
  3 \leftrightarrow \mathrm{U}.
\]
The variable $x_i$ has cardinality 4 for all $i$. The EDA population is an
$N \times n$ integer matrix, and all probabilistic models operate over this
quaternary alphabet.

\subsection{Objective Functions}

The \texttt{rna\_design} project implements three complementary scalar
objective functions. Each measures a distinct aspect of how well a designed
sequence $\mathbf{x}$ achieves the target structure $s^*$. All three are
formulated as \emph{maximisation} objectives by negating their natural
cost.

% ---- Objective 1 -----------------------------------------------------------
\subsubsection{Objective 1: Negative Base-Pair Distance}
\label{sec:obj1}

Let $s(\mathbf{x}) = \mathcal{F}(\mathbf{x})$ be the MFE secondary structure
predicted for sequence $\mathbf{x}$ by the Vienna RNA package~\cite{lorenz2011viennarna}. The
\emph{base-pair distance} between $s(\mathbf{x})$ and the target $s^*$ is
\[
  d_{\mathrm{BP}}(\mathbf{x},\, s^*) \;=\;
    \bigl|\mathcal{B}(s(\mathbf{x})) \,\triangle\, \mathcal{B}(s^*)\bigr|,
\]
where $\mathcal{B}(s)$ is the set of base pairs encoded by structure $s$ and
$\triangle$ denotes the symmetric difference. Intuitively,
$d_{\mathrm{BP}}$ counts the number of base pairs that are present in one
structure but absent in the other.

The \textbf{negative base-pair distance} is used as the first objective:
\[
  f_1(\mathbf{x}) = -\,d_{\mathrm{BP}}(\mathbf{x},\, s^*).
\]
A solution for which $f_1(\mathbf{x}) = 0$ has solved the design problem
exactly: the predicted MFE structure matches the target perfectly.

\paragraph{Benefits.} $f_1$ is the most direct measure of design quality and
corresponds to the standard success criterion used in benchmarks such as
Eterna100~\cite{lee2014eterna}. It is efficient to compute since it requires
only one call to an MFE folding routine.

\paragraph{Limitations.} The MFE structure is a single point prediction and
does not account for the thermodynamic ensemble of structures that the
sequence samples at physiological temperature. A sequence can achieve
$f_1(\mathbf{x}) = 0$ while still spending significant probability on
alternative, non-target folds.

% ---- Objective 2 -----------------------------------------------------------
\subsubsection{Objective 2: Probability of the Target Structure}
\label{sec:obj2}

Given a sequence $\mathbf{x}$, the Boltzmann partition function $Z(\mathbf{x})$ sums
the Boltzmann weights of all possible secondary structures:
\[
  Z(\mathbf{x}) = \sum_{s} \exp\!\left(-\frac{\Delta G(s, \mathbf{x})}{RT}\right),
\]
where $\Delta G(s, \mathbf{x})$ is the free energy of structure $s$ folded
from sequence $\mathbf{x}$, $R$ is the gas constant, and $T$ is the absolute
temperature. The probability of the target structure $s^*$ under the
Boltzmann ensemble is
\[
  p(s^* \mid \mathbf{x}) = \frac{\exp\!\left(-\Delta G(s^*, \mathbf{x})\,/\,RT\right)}{Z(\mathbf{x})}.
\]
The second objective is the log-probability of the target:
\[
  f_2(\mathbf{x}) = \log p(s^* \mid \mathbf{x})
               = -\frac{\Delta G(s^*, \mathbf{x})}{RT} - \log Z(\mathbf{x}).
\]
Maximising $f_2$ pushes the sequence toward assigning high Boltzmann weight to
the target structure relative to all competing folds.

\paragraph{Benefits.} $f_2$ is a principled thermodynamic measure that
accounts for the full ensemble, not just the MFE structure. Sequences with
high $f_2$ are more likely to adopt the target structure under realistic
folding conditions.

\paragraph{Limitations.} Computing the partition function $Z(\mathbf{x})$ via
the McCaskill algorithm~\cite{mccaskill1990equilibrium} has time complexity
$O(n^3)$ and is thus considerably more expensive than MFE folding. For long
sequences or large populations this cost may be prohibitive.

% ---- Objective 3 -----------------------------------------------------------
\subsubsection{Objective 3: Negative Ensemble Defect}
\label{sec:obj3}

The \emph{ensemble defect}~\cite{zadeh2011nupack} quantifies the expected
number of incorrectly paired nucleotides when a structure is drawn at random
from the Boltzmann ensemble:
\[
  D(\mathbf{x},\, s^*) = n \;-\;
    \sum_{i=1}^{n}\sum_{j=0}^{n} \mathbb{1}[s^*_{ij}]\, p_{ij}(\mathbf{x}),
\]
where $p_{ij}(\mathbf{x})$ is the base-pair probability between positions $i$
and $j$ (with $j=0$ denoting the unpaired state), and $\mathbb{1}[s^*_{ij}]$
is 1 if positions $i$ and $j$ are paired (or both unpaired) in $s^*$. The
\emph{normalised ensemble defect} divides by $n$:
\[
  \tilde{D}(\mathbf{x},\, s^*) = \frac{D(\mathbf{x},\, s^*)}{n} \in [0,1].
\]
The third objective is the negative normalised ensemble defect:
\[
  f_3(\mathbf{x}) = -\,\tilde{D}(\mathbf{x},\, s^*).
\]
A value of $f_3(\mathbf{x}) = 0$ means that, on average, every nucleotide
adopts its target pairing state across the Boltzmann ensemble.

\paragraph{Benefits.} The ensemble defect is a position-level metric:
it not only penalises wrong base pairs globally but identifies which
specific positions are least likely to occupy their target pairing state.
It is therefore a more sensitive guide for the search than $f_1$, especially
in the early phases of optimisation.

\paragraph{Limitations.} Like $f_2$, computing $\tilde{D}$ requires the full
base-pair probability matrix, incurring $O(n^3)$ cost. Furthermore, $f_3$
can be zero for sequences that nonetheless have $f_1 < 0$ (i.e., the average
is correct but the most probable single structure differs from $s^*$).

% ─────────────────────────────────────────────────────────────────────────────
\section{Description of the EDA Approaches}
\label{sec:edas}
% ─────────────────────────────────────────────────────────────────────────────

Estimation of Distribution Algorithms (EDAs)~\cite{larranaga2002eda,
pelikan2002eda} replace the crossover and mutation operators of classical
evolutionary algorithms with an explicit probabilistic model of the selected
population. The generic EDA loop is shown in Algorithm~\ref{alg:eda}.

\begin{algorithm}[h]
\caption{Generic EDA}
\label{alg:eda}
\begin{algorithmic}[1]
\State $t \leftarrow 0$. Initialise population $D_0$ with $N$ random solutions.
\Repeat
  \State Evaluate $D_t$ using the fitness function.
  \State Select $D_t^S \subseteq D_t$ of $K \leq N$ promising solutions.
  \State Learn probabilistic model $\hat{p}_t(\mathbf{x})$ from $D_t^S$.
  \State Sample $D_{t+1}$ from $\hat{p}_t(\mathbf{x})$.
  \State $t \leftarrow t + 1$.
\Until{stopping criterion is met}
\end{algorithmic}
\end{algorithm}

The \texttt{rna\_design} project uses four EDA variants that differ in the
class of probabilistic model learned in step~5. All implementations reside
in the \texttt{pateda} library (module \texttt{pateda.learning}).

\subsection{UMDA — Univariate Marginal Distribution Algorithm}

UMDA~\cite{muhlenbein1997umda} assumes complete statistical independence
among the $n$ variables. The model factorises as a product of univariate
marginals:
\[
  \hat{p}_t(\mathbf{x}) = \prod_{i=1}^{n} p_t(x_i),
\]
where each $p_t(x_i)$ is estimated by frequency counting over the selected
population $D_t^S$:
\[
  p_t(x_i = v) = \frac{|\{s \in D_t^S : s_i = v\}|}{|D_t^S|},
  \quad v \in \{0,1,2,3\}.
\]
Laplace smoothing with parameter $\alpha \geq 0$ prevents zero probabilities:
\[
  \hat{p}_t(x_i = v) = \frac{\text{count}(x_i = v) + \alpha}{|D_t^S| + 4\alpha}.
\]
Sampling from the UMDA model is $O(nN)$; learning is $O(nK)$. The model
contains $4n$ free parameters (four probabilities per position), making it
by far the lightest of the four variants.

\paragraph{Strengths and limitations.} UMDA is computationally cheap and
robust to small population sizes. However, the independence assumption is
strongly violated in RNA design: base pairs by definition couple two
positions, and the nucleotide at a paired position is severely constrained
by the complementarity rules (A–U, G–C, G–U). Ignoring these dependencies
slows convergence and often yields suboptimal solutions.

\subsection{Tree-EDA}

Tree-EDA~\cite{baluja1997dependency, santana2008protein} relaxes the
independence assumption by allowing each variable to depend on \emph{at
most one} other variable (its parent), forming a directed forest. The
distribution factorises as
\[
  \hat{p}_t(\mathbf{x})
  = p_t(x_\mathrm{root}) \prod_{i \neq \mathrm{root}} p_t(x_i \mid x_{\pi(i)}),
\]
where $\pi(i)$ is the unique parent of node $i$ in the tree.

\subsubsection*{Structure Learning (Chow--Liu algorithm)}

\begin{enumerate}
  \item Compute the pairwise mutual information for all $\binom{n}{2}$ pairs:
    \[
      I(X_i;\, X_j) = \sum_{v,w} p(x_i{=}v,\, x_j{=}w)
        \log \frac{p(x_i{=}v,\, x_j{=}w)}{p(x_i{=}v)\,p(x_j{=}w)}.
    \]
  \item Construct the complete weighted graph $\mathcal{G}$ on $n$ nodes
    with edge weight $I(X_i; X_j)$.
  \item Find the \emph{maximum weight spanning tree} (MWST) of $\mathcal{G}$
    using Prim's algorithm. This is the tree structure that best approximates
    the joint distribution in KL-divergence~\cite{chow1968tree}.
  \item Root the tree at the node involved in the highest-weight edge; orient
    all edges away from the root.
\end{enumerate}

\subsubsection*{Parameter learning}

For the root: $p_t(x_r)$ is estimated from univariate counts. For each
non-root node $i$: $p_t(x_i \mid x_{\pi(i)})$ is estimated from bivariate
counts over $D_t^S$.

\paragraph{Complexity.} Learning the MWST requires computing $O(n^2)$
pairwise mutual information values and running Prim's algorithm in
$O(n^2)$ time. Sampling is $O(nN)$ via probabilistic logic sampling (PLS).

\paragraph{Strengths and limitations.} Tree-EDA can capture pairwise
dependencies — exactly the dependencies induced by base pairs — making it
substantially more powerful than UMDA for RNA design. Its main limitation
is that the MWST must connect \emph{all} $n$ nodes: if only $m \ll
\binom{n}{2}$ pairs truly interact, the algorithm still wastes time
computing mutual information for the remaining $\binom{n}{2} - m$ pairs.
Furthermore, tree edges are limited to one parent per node, so higher-order
interactions in stacked base pairs or pseudoknot-like structures cannot be
fully captured.

\subsection{Tree-EDA\textsubscript{r} — Restricted Tree EDA}

Tree-EDA\textsubscript{r}~\cite{santana2012transfer} is a variant of
Tree-EDA that incorporates prior knowledge about which pairs of variables
\emph{can} interact. This knowledge is encoded in a binary symmetric
\emph{interaction matrix} $R \in \{0,1\}^{n \times n}$, where $R_{ij} = 1$
means that variables $X_i$ and $X_j$ are structurally related and their
mutual information should be considered for the spanning tree.

\subsubsection*{Modified Structure Learning}

Step~1 of the Chow--Liu algorithm is restricted so that mutual information
is computed \emph{only} for interacting pairs:
\[
  \tilde{I}(X_i;\, X_j) =
  \begin{cases}
    I(X_i;\, X_j) & \text{if } R_{ij} = 1,\\
    0              & \text{if } R_{ij} = 0.
  \end{cases}
\]
The MWST is then computed on the sparse graph induced by $R$. Non-interacting
pairs have weight 0 and therefore cannot appear as tree edges.

The implementation in \texttt{pateda/learning/tree\_r.py} (class
\texttt{LearnTreeModelR}) accepts \texttt{interaction\_matrix} at
construction time and replaces the full mutual information loop with a loop
over interacting pairs only:

\begin{center}
\begin{minipage}{0.85\textwidth}
\begin{verbatim}
for i in range(n_vars - 1):
    for j in range(i + 1, n_vars):
        if interaction_matrix[i, j] == 0:
            continue          # skip non-interacting pair
        mi = compute_MI(i, j, ...)
        mi_matrix[i, j] = mi / (card_i * card_j)
        mi_matrix[j, i] = mi_matrix[i, j]
\end{verbatim}
\end{minipage}
\end{center}

\paragraph{Strengths.} When $m = \sum_{i<j} R_{ij} \ll \binom{n}{2}$,
Tree-EDA\textsubscript{r} computes only $m$ mutual information values instead
of $\binom{n}{2}$, yielding a speedup proportional to $\binom{n}{2}/m$.
More importantly, restricting the MWST to structurally relevant edges
prevents the algorithm from fitting spurious statistical dependencies that
arise in finite populations, which can mislead the search.

\paragraph{Limitations.} The quality of Tree-EDA\textsubscript{r} depends
entirely on the quality of the interaction matrix $R$. If $R$ omits true
interactions, the spanning tree will be forced to ignore them. The method is
therefore most effective when the problem structure is well-characterised a
priori — which is precisely the case in RNA design, where base pairing is
known from the target structure.

\subsection{MK-EDA1 — First-Order Markov Chain EDA}

MK-EDA1 is the first-order ($k=1$) instance of the $k$-order Markov Chain
EDA~\cite{santana2008protein} implemented in
\texttt{pateda/learning/markov.py} (\texttt{LearnMarkovChain(k=1)}). The
probabilistic model factorises as
\[
  \hat{p}_t(\mathbf{x})
  = p_t(x_1,\, x_2)
    \prod_{i=3}^{n} p_t(x_i \mid x_{i-1}),
\]
where the joint distribution over the first two variables and all pairwise
conditional distributions are estimated from $D_t^S$.

\subsubsection*{Structure}

The model forms a fixed \emph{chain} over the natural left-to-right ordering
of the sequence. The clique structure is:
\[
  \{x_1, x_2\},\quad \{x_2, x_3\},\quad \ldots,\quad \{x_{n-1}, x_n\}.
\]
Each clique $\{x_{i-1}, x_i\}$ has one overlapping variable ($x_{i-1}$,
already instantiated) and one new variable ($x_i$). Sampling proceeds
left-to-right by ancestral sampling: draw $x_1, x_2$ jointly, then draw
each $x_i$ conditioned on the already-sampled $x_{i-1}$.

\subsubsection*{Parameter learning}

\begin{itemize}
  \item \textbf{Initial clique} $\{x_1, x_2\}$: joint counts over
    $4 \times 4 = 16$ configurations, normalised.
  \item \textbf{Transition cliques} $\{x_{i-1}, x_i\}$ for $i \geq 3$:
    conditional probability tables $p_t(x_i = w \mid x_{i-1} = v)$ estimated
    from frequency counts, with Laplace smoothing $\alpha = 0.1$.
\end{itemize}

\paragraph{Strengths.} MK-EDA1 captures \emph{consecutive} pairwise
dependencies and is natural for RNA sequences, where positions adjacent in
the chain often co-vary to maintain local secondary structure elements such
as hairpin loops and bulges. The number of cliques is $n-1$, giving linear
$O(n)$ learning and sampling complexity.

\paragraph{Limitations.} MK-EDA1 is blind to \emph{long-range} base-pair
interactions, which are ubiquitous in RNA secondary structure. A stem formed
by base pairs $(i, j)$ with $|i - j| \gg 1$ is entirely invisible to the
chain model. Furthermore, the chain ordering imposes a fixed dependency
direction, so MK-EDA1 cannot selectively emphasise known interacting pairs
the way Tree-EDA\textsubscript{r} does.

\subsection{Summary of Differences}

Table~\ref{tab:edas} contrasts the four EDA variants across the dimensions
most relevant to RNA design.

\begin{table}[h]
\centering
\caption{Comparison of EDA variants used in \texttt{rna\_design}.}
\label{tab:edas}
\begin{tabular}{lcccc}
\toprule
\textbf{Property} & \textbf{UMDA} & \textbf{Tree-EDA}
  & \textbf{Tree-EDA\textsubscript{r}} & \textbf{MK-EDA1} \\
\midrule
Model family           & Univariate    & Tree (BN)  & Tree (BN)     & Markov chain \\
Max. dependency order  & 0             & 1          & 1             & 1            \\
Dependency structure   & None          & Learned    & Learned+Prior & Fixed chain  \\
Long-range pairs       & \texttimes    & \checkmark & \checkmark    & \texttimes   \\
Uses prior $R$         & \texttimes    & \texttimes & \checkmark    & \texttimes   \\
MI pairs computed      & 0             & $n(n-1)/2$ & $\leq n(n-1)/2$ & $n-1$      \\
Learning complexity    & $O(nK)$       & $O(n^2 K)$ & $O(m K)$      & $O(nK)$     \\
Parameters             & $4n$          & $O(16n)$   & $O(16n)$      & $O(16n)$    \\
\bottomrule
\end{tabular}
\end{table}
\noindent Here $K$ is the size of the selected population and
$m = \sum_{i<j} R_{ij}$ is the number of interacting pairs in the
interaction matrix (with $m \leq n(n-1)/2$).

% ─────────────────────────────────────────────────────────────────────────────
\section{Using Problem Interactions to Bias the Search}
\label{sec:interactions}
% ─────────────────────────────────────────────────────────────────────────────

\subsection{Motivation}

The RNA design problem belongs to the class of \emph{gray-box optimisation}
problems~\cite{santana2017graybox}: the target structure is known a priori
and encodes, in explicit form, the structural relationships among the
variables. Specifically, a base pair $(i, j)$ in $s^*$ tells us unambiguously
that nucleotides at positions $i$ and $j$ are strongly coupled: the identity
of $x_i$ severely constrains the set of admissible values for $x_j$
(complementarity rules). Moreover, consecutive positions share stacking
interactions and backbone geometry constraints. Both types of information
are available before the EDA runs a single generation.

Tree-EDA\textsubscript{r} is designed to exploit this gray-box information:
instead of discovering interactions from scratch via mutual information, it
\emph{restricts} the set of candidate edges in the MWST to those variable
pairs that are structurally related according to the target structure.

\subsection{The \texttt{find\_matrix\_interactions\_RNA\_design} Function}

The function \texttt{find\_matrix\_interactions\_RNA\_design(RNA\_structure)},
implemented in \texttt{external/pateda/learning/interaction\_learning.py},
constructs the binary interaction matrix $R$ directly from the Dot-Bracket
string $s^*$ of length $n$. It encodes two complementary types of structural
interaction.

\subsubsection*{Type 1 — Sequence Adjacency}

Consecutive nucleotides in the RNA backbone are covalently bonded.
Their identities jointly influence local folding stability through
base-stacking interactions, which account for a significant fraction of
the total free energy of RNA secondary structures~\cite{mathews2004thermodynamics}. Accordingly,
\[
  R_{i,\,i+1} = R_{i+1,\,i} = 1 \quad \forall\, i \in \{1, \ldots, n-1\}.
\]
This produces the $n-1$ edges of a path graph on the $n$ positions.

\subsubsection*{Type 2 — Base-Pair Interactions}

A base pair $(i, j)$ in $s^*$ (with $i < j$) is identified by a matched
\texttt{(}\,/\,\texttt{)} in the Dot-Bracket string. The function uses
a stack-based parser to recover all such pairs:
\begin{center}
\begin{minipage}{0.85\textwidth}
\begin{verbatim}
stack = []
for position i, symbol in enumerate(RNA_structure):
    if symbol == '(':
        stack.append(i)
    elif symbol == ')':
        j = stack.pop()       # j is the matching '(' position
        R[j, i] = R[i, j] = 1
\end{verbatim}
\end{minipage}
\end{center}
For each matched pair $(j, i)$ discovered this way, the matrix is set
symmetrically:
\[
  R_{i,j} = R_{j,i} = 1.
\]
Unpaired positions (\texttt{.}) contribute no base-pair edges; they are
covered only by sequence-adjacency edges with their two immediate neighbours.

\subsubsection*{Properties of $R$}

Let $B$ denote the number of base pairs encoded by $s^*$. The matrix $R$
satisfies:
\begin{itemize}
  \item \textbf{Binary}: $R_{ij} \in \{0,1\}$ for all $i, j$.
  \item \textbf{Symmetric}: $R_{ij} = R_{ji}$.
  \item \textbf{Zero diagonal}: $R_{ii} = 0$.
  \item \textbf{Interaction count}: exactly $m = (n-1) + B$ edges are set
    to 1 in the upper triangle of $R$, minus any overlap where a base pair
    happens to link adjacent positions (which is rare but possible in
    degenerate cases).
\end{itemize}

For a typical RNA secondary structure, $B \approx n/2$ (roughly half the
positions are paired), giving $m \approx 3n/2$. Compared to the $\binom{n}{2}
\approx n^2/2$ pairs that Tree-EDA would test, this represents a speedup of
order $n/3$ in the mutual-information computation step — highly significant
for long sequences.

\subsection{Passing $R$ to Tree-EDA\textsubscript{r}}

The typical usage pattern in \texttt{rna\_design} is:
\begin{center}
\begin{minipage}{0.85\textwidth}
\begin{verbatim}
from pateda.learning.interaction_learning import \
    find_matrix_interactions_RNA_design
from pateda.learning.tree_r import LearnTreeModelR

R = find_matrix_interactions_RNA_design(target_structure)
learner = LearnTreeModelR(interaction_matrix=R, alpha=0.0)
\end{verbatim}
\end{minipage}
\end{center}
At each generation, \texttt{LearnTreeModelR.learn()} computes mutual
information only for the $m$ pairs with $R_{ij}=1$, then runs Prim's
MWST algorithm on the resulting sparse graph. The resulting tree is
guaranteed to contain only edges that correspond to either backbone
adjacency or Watson--Crick base pairing in the target — the two
dominant sources of epistasis in RNA sequences.

% ─────────────────────────────────────────────────────────────────────────────
\section{Description of the Eterna100 Database}
\label{sec:eterna}
% ─────────────────────────────────────────────────────────────────────────────

\subsection{Overview}

The \textbf{Eterna100}~\cite{lee2014eterna, anderson2016eterna} is a widely
adopted benchmark suite for RNA design algorithms. It consists of
\textbf{100 RNA secondary structure puzzles} curated from the EteRNA
citizen-science game~\cite{das2010eterna}, in which human players design
RNA sequences to fold into prescribed target structures. The benchmark was
introduced specifically to evaluate computational design algorithms under
conditions that reflect real-world difficulty.

\subsection{Construction and Curation}

The 100 puzzles were selected from the EteRNA game database to cover a broad
range of structural complexity. Selection criteria included:
\begin{itemize}
  \item \textbf{Difficulty}: puzzles solved by fewer human players are harder
    and thus more discriminating for algorithmic comparison.
  \item \textbf{Diversity}: puzzles span a wide range of sequence lengths
    (from approximately 14 to over 400 nucleotides) and structural motifs
    (hairpin loops, internal loops, bulges, multi-loops, and pseudoknot-free
    structures of varying nesting depth).
  \item \textbf{Validity}: all target structures are thermodynamically
    realisable under the Turner energy model used by Vienna RNA.
\end{itemize}

\subsection{Evaluation Protocol}

The standard Eterna100 evaluation protocol, as used in \texttt{rna\_design},
is as follows:
\begin{itemize}
  \item An algorithm is given the target structure $s^*_k$ for puzzle
    $k \in \{1, \ldots, 100\}$ and a computational budget (typically
    measured in number of objective function evaluations or wall-clock time).
  \item The algorithm returns a candidate sequence $\hat{\mathbf{x}}_k$.
  \item The sequence is folded with the Vienna RNA package
    (\texttt{RNAfold}); if the predicted MFE structure equals $s^*_k$
    (i.e., base-pair distance $= 0$), the puzzle is counted as
    \emph{solved}.
  \item The primary performance metric is the \textbf{solve rate}: the
    number of puzzles solved out of 100.
\end{itemize}
Secondary metrics include the average base-pair distance for unsolved
puzzles and the average number of evaluations required to find a solution.

\subsection{Structural Characteristics}

Table~\ref{tab:eterna} summarises the main structural statistics of the
Eterna100 benchmark as reported in~\cite{lee2014eterna, andronescu2016eterna}.

\begin{table}[h]
\centering
\caption{Structural statistics of the Eterna100 benchmark.}
\label{tab:eterna}
\begin{tabular}{lcc}
\toprule
\textbf{Property} & \textbf{Min} & \textbf{Max} \\
\midrule
Sequence length $n$          & 14   & 400+  \\
Number of base pairs $B$     & 4    & 160+  \\
Nesting depth                & 1    & 8+    \\
Number of multi-loops        & 0    & 10+   \\
Solve rate (top human)       & ---  & 100\% \\
\bottomrule
\end{tabular}
\end{table}

\subsection{Relevance to EDA Benchmarking}

The Eterna100 benchmark is particularly well-suited to evaluating EDA
approaches because:
\begin{enumerate}
  \item \textbf{Problem diversity}: the 100 puzzles span a continuum from
    trivial (short hairpins) to very hard (long multi-loop structures),
    enabling fine-grained analysis of algorithm scaling.
  \item \textbf{Known interactions}: the target structure provides an
    explicit interaction matrix $R$ via
    \texttt{find\_matrix\_interactions\_RNA\_design}, enabling a controlled
    comparison between unrestricted and restricted EDAs.
  \item \textbf{Standard baseline}: solve rates from many alternative
    algorithms (RNAINVERSE, INFO-RNA, antaRNA, MCTS-RNA, deep-learning
    methods) are available in the literature, permitting direct comparison.
\end{enumerate}

% ─────────────────────────────────────────────────────────────────────────────
\section{Related Approaches}
\label{sec:related}
% ─────────────────────────────────────────────────────────────────────────────

The \texttt{other\_rna\_implementations/} directory collects reference
implementations of established RNA design algorithms. These implementations
serve two roles in \texttt{rna\_design}: (i) they provide the
thermodynamic objective function computations (base-pair distance, ensemble
defect, partition function) that are wrapped into the objective functions
$f_1$, $f_2$, $f_3$ described in Section~\ref{sec:representation}, and (ii)
they constitute the baseline algorithms against which EDA performance is
compared.

\subsection{RNAINVERSE (Vienna RNA)}

\textbf{RNAINVERSE}~\cite{hofacker1994fast, lorenz2011viennarna} is the
classical RNA design tool distributed as part of the Vienna RNA package. It
is a stochastic local-search algorithm that uses a mutation-and-acceptance
loop over the space of sequences:
\begin{enumerate}
  \item Initialise $\mathbf{x}$ randomly, respecting Watson--Crick
    complementarity at base-paired positions.
  \item Repeat: mutate a randomly chosen position; accept the mutant if
    it reduces the Hamming distance between the predicted MFE structure
    and $s^*$.
\end{enumerate}
\textbf{Commonalities with \texttt{rna\_design}}: RNAINVERSE uses the same
Vienna RNA thermodynamic model for MFE prediction; \texttt{rna\_design}
wraps the \texttt{RNAfold} and \texttt{RNAdistance} utilities from the
same package to compute $f_1$, $f_2$, and $f_3$.

\textbf{Differences}: RNAINVERSE is a single-solution local-search method with
no population and no probabilistic model. It cannot exploit global
statistical patterns in the search history, is prone to local optima, and
typically requires multiple restarts to solve hard puzzles. In contrast, the
EDAs in \texttt{rna\_design} maintain a population and build an explicit
model of the promising sequence regions, enabling a more global and
information-guided search.

\subsection{INFO-RNA}

\textbf{INFO-RNA}~\cite{busch2006informa} extends the local-search idea by
incorporating a preprocessing step that initialises the sequence using
a dynamic programming (DP) procedure. The DP computes, for each position,
the most probable nucleotide given its structural context in $s^*$:
paired positions are initialised with the most frequent base-pair type
in a reference dataset; unpaired positions with the most frequent unpaired
nucleotide.

\textbf{Commonalities with \texttt{rna\_design}}: INFO-RNA's initialisation
heuristic is conceptually related to the interaction-aware seeding that
\texttt{rna\_design} can apply when using $R$ from
\texttt{find\_matrix\_interactions\_RNA\_design}: both methods exploit
the known pairing structure of $s^*$ to bias the initial sequence.

\textbf{Differences}: INFO-RNA combines DP initialisation with a local
stochastic search that is still single-solution and restart-dependent. The
EDAs in \texttt{rna\_design} build a full probabilistic model over all $n$
positions simultaneously and use population-level statistics to escape local
optima.

\subsection{antaRNA}

\textbf{antaRNA}~\cite{doelker2015antarna} (ant colony optimisation for RNA
design) models the sequence-design problem as a graph traversal and uses
simulated-annealing-based ant colony optimization. Ants deposit pheromones
on nucleotide choices that produce structures closer to the target. Additional
constraints such as GC content and sequence motifs can be incorporated as
soft constraints.

\textbf{Commonalities with \texttt{rna\_design}}: antaRNA uses a population
of agents (ants) and maintains marginal statistics over nucleotide choices,
resembling the univariate model of UMDA. Like Tree-EDA\textsubscript{r},
antaRNA can incorporate a priori information about the structure to prune the
search space.

\textbf{Differences}: the pheromone update rule in antaRNA is position-wise
(analogous to UMDA) and does not model pairwise or higher-order dependencies.
The EDAs in \texttt{rna\_design} go further by capturing pairwise
co-evolution between base-paired positions, which antaRNA approximates only
implicitly through the constraint that paired positions must be
complementary.

\subsection{MCTS-RNA}

\textbf{MCTS-RNA}~\cite{yang2017mcts} applies Monte Carlo Tree Search to RNA
design. Sequence positions are filled left-to-right; at each step, MCTS
simulates random completions of the partial sequence, estimates the value
of each nucleotide choice, and selects the most promising option.

\textbf{Commonalities with \texttt{rna\_design}}: MCTS-RNA's sequential
position-by-position construction order is similar to the left-to-right
ancestral sampling in MK-EDA1. Both methods condition nucleotide choices on
previously assigned positions and use rollout or model information to guide
the decision.

\textbf{Differences}: MCTS-RNA is a tree-search method (not a population-based
EDA) and its rollout policy does not maintain a generative probability model
over the full sequence. The EDAs in \texttt{rna\_design} learn and re-use an
explicit generative model at every generation, enabling the progressive
sharpening of marginal and conditional distributions as the search advances.

\subsection{Relationship Summary}

Table~\ref{tab:related} places the related methods and the four EDA variants
in a common framework.

\begin{table}[h]
\centering
\caption{Comparison of \texttt{rna\_design} EDA variants and related algorithms.}
\label{tab:related}
\begin{tabular}{lccccc}
\toprule
\textbf{Algorithm} & \textbf{Population} & \textbf{Model} & \textbf{Pairwise deps.}
  & \textbf{Structural prior} & \textbf{Source} \\
\midrule
RNAINVERSE   & \texttimes & None               & \texttimes & Implicit & Vienna RNA \\
INFO-RNA      & \texttimes & DP init            & \texttimes & \checkmark & other\_rna \\
antaRNA       & \checkmark & Pheromone (univ.)  & \texttimes & Soft       & other\_rna \\
MCTS-RNA      & \texttimes & Rollout            & Partial    & \texttimes & other\_rna \\
\midrule
UMDA          & \checkmark & Univariate         & \texttimes & \texttimes & pateda \\
MK-EDA1       & \checkmark & Markov chain       & Adjacent   & \texttimes & pateda \\
Tree-EDA      & \checkmark & Tree (BN)          & \checkmark & \texttimes & pateda \\
Tree-EDA\textsubscript{r} & \checkmark & Tree (BN) & \checkmark & \checkmark & pateda \\
\bottomrule
\end{tabular}
\end{table}

The key differentiator of Tree-EDA\textsubscript{r} within this landscape is
the combination of a population-based search, a pairwise probabilistic model,
and explicit incorporation of the structural prior $R$ derived from the target
structure. This combination aligns it most closely with the gray-box
optimisation paradigm~\cite{santana2017graybox}: it does not treat the
objective function as a black box but instead uses available structural
knowledge to focus the model-learning step on the interactions that matter
most for RNA design.

% ─────────────────────────────────────────────────────────────────────────────
\bibliographystyle{plain}
\begin{thebibliography}{99}

\bibitem{tinoco1999rna}
I. Tinoco and C. Bustamante,
``How RNA folds,''
\emph{Journal of Molecular Biology}, 293(2):271--281, 1999.

\bibitem{zuker1989mfold}
M. Zuker and P. Stiegler,
``Optimal computer folding of large RNA sequences using thermodynamics and auxiliary information,''
\emph{Nucleic Acids Research}, 9(1):133--148, 1981.

\bibitem{hofacker1994fast}
I.~L. Hofacker, W. Fontana, P.~F. Stadler, L.~S. Bonhoeffer, M. Tacker, and P. Schuster,
``Fast folding and comparison of RNA secondary structures,''
\emph{Monatshefte f\"ur Chemie}, 125(2):167--188, 1994.

\bibitem{lorenz2011viennarna}
R. Lorenz, S.~H. Bernhart, C.~H. Zu Siederdissen, H. Tafer, C. Flamm, P.~F. Stadler, and I.~L. Hofacker,
``{ViennaRNA} package 2.0,''
\emph{Algorithms for Molecular Biology}, 6(1):26, 2011.

\bibitem{mccaskill1990equilibrium}
J.~S. McCaskill,
``The equilibrium partition function and base pair binding probabilities for RNA secondary structure,''
\emph{Biopolymers}, 29(6-7):1105--1119, 1990.

\bibitem{zadeh2011nupack}
J.~N. Zadeh, C.~D. Steenberg, J.~S. Bois, B.~R. Wolfe, M.~B. Pierce, A.~R. Khan, R.~M. Dirks, and N.~A. Pierce,
``{NUPACK}: Analysis and design of nucleic acid systems,''
\emph{Journal of Computational Chemistry}, 32(1):170--173, 2011.

\bibitem{wachsmuth2013riboswitch}
M. Wachsmuth, S. Findeiss, N. Weissheimer, P.~F. Stadler, and M. M\"orl,
``De novo design of a synthetic riboswitch that regulates transcription termination,''
\emph{Nucleic Acids Research}, 41(4):2541--2551, 2013.

\bibitem{lee2014eterna}
J. Lee, H. Kladwang, M. Lee, D. Cantu, M. Azizyan, H. Kim, A. Limpaecher, S. Gatta, S. Testa, A. Treuille, and R. Das,
``RNA design rules from a massive open laboratory,''
\emph{Proceedings of the National Academy of Sciences}, 111(6):2122--2127, 2014.

\bibitem{anderson2016eterna}
E. Anderson and R. Das,
``{Eterna100} benchmark for RNA inverse folding,''
available at \texttt{https://github.com/eternagame/Eterna100-benchmarking}, 2016.

\bibitem{andronescu2016eterna}
M. Andronescu, A. Condon, H.~H. Hoos, D.~H. Mathews, and K.~P. Murphy,
``Efficient parameter estimation for RNA secondary structure prediction,''
\emph{Bioinformatics}, 23(13):i19--i28, 2007.

\bibitem{mathews2004thermodynamics}
D.~H. Mathews and D.~H. Turner,
``Prediction of RNA secondary structure by free energy minimization,''
\emph{Current Opinion in Structural Biology}, 16(3):270--278, 2006.

\bibitem{larranaga2002eda}
P. Larra\~naga and J.~A. Lozano (Eds.),
\emph{Estimation of Distribution Algorithms: A New Tool for Evolutionary Computation}.
Kluwer Academic Publishers, 2002.

\bibitem{pelikan2002eda}
M. Pelikan, D.~E. Goldberg, and F.~G. Lobo,
``A survey of optimization by building and using probabilistic models,''
\emph{Computational Optimization and Applications}, 21(1):5--20, 2002.

\bibitem{muhlenbein1997umda}
H. M\"uhlenbein and G. Paa\ss,
``From recombination of genes to the estimation of distributions~I: Binary parameters,''
in \emph{Parallel Problem Solving from Nature (PPSN-IV)}, LNCS 1411, pp.~178--187, 1996.

\bibitem{baluja1997dependency}
S. Baluja and S. Davies,
``Using optimal dependency-trees for combinatorial optimization: Learning the structure of the search space,''
in \emph{Proceedings of ICML-1997}, pp.~30--38, 1997.

\bibitem{chow1968tree}
C. Chow and C. Liu,
``Approximating discrete probability distributions with dependence trees,''
\emph{IEEE Transactions on Information Theory}, 14(3):462--467, 1968.

\bibitem{santana2008protein}
R. Santana, P. Larra\~naga, and J.~A. Lozano,
``Protein folding in simplified models with estimation of distribution algorithms,''
\emph{IEEE Transactions on Evolutionary Computation}, 12(4):418--438, 2008.

\bibitem{santana2012transfer}
R. Santana, A. Mendiburu, and J.~A. Lozano,
``Structural transfer using EDAs: An application to multi-marker tagging SNP selection,''
in \emph{Proceedings of the 2012 IEEE Congress on Evolutionary Computation (CEC-2012)},
pp.~3484--3491, Brisbane, Australia, 2012.

\bibitem{santana2017graybox}
R. Santana,
``Gray-box optimization and factorized distribution algorithms: where two worlds collide,''
\emph{arXiv preprint arXiv:1707.03093}, 2017.

\bibitem{busch2006informa}
A. Busch and R. Backofen,
``{INFO-RNA} --- a fast approach to inverse RNA folding,''
\emph{Bioinformatics}, 22(15):1823--1831, 2006.

\bibitem{doelker2015antarna}
R. D\"oelker, M. Mann, and R. Backofen,
``{antaRNA}: Ant colony-based RNA sequence design,''
\emph{Bioinformatics}, 31(19):3114--3121, 2015.

\bibitem{yang2017mcts}
X. Yang, K. Yoshizoe, A. Taneda, and K. Tsuda,
``{RNA} inverse folding using Monte Carlo tree search,''
\emph{BMC Bioinformatics}, 18(1):468, 2017.

\bibitem{das2010eterna}
R. Das, J. Karanicolas, and D. Baker,
``Atomic accuracy in predicting and designing noncanonical RNA structure,''
\emph{Nature Methods}, 7(4):291--294, 2010.

\end{thebibliography}

\end{document}
